% Document/LinearAlgebra/VectorSpaces/Matrices/ChangeOfBasis.tex
% Phase 17: Change of Basis Matrices; Inj/Surj/Bij via Matrices; Invertibility

\newpage
\section{Change of Basis}
\label{sec:change-of-basis}

\subsection{Change of Basis Matrices}

\begin{definition}[Change of Basis Matrix]
\label{def:cob-matrix}
    Let $V$ be a $K$-vector space with two ordered bases $\mathcal{B}$ and $\mathcal{C}$.
    The \textbf{change of basis matrix from $\mathcal{B}$ to $\mathcal{C}$} is:
    \[
        \CoordMatrix[\mathcal{B}][\mathcal{C}]{\Identity{V}}
    \]
    Its $j$-th column is $\CoordVec[\mathcal{C}]{b_j}$ --- the coordinate vector of the
    $j$-th basis vector of $\mathcal{B}$ expressed in basis $\mathcal{C}$.
\end{definition}

\begin{proposition}[Properties of Change of Basis Matrices]
\label{prop:cob-properties}
    Let $V$ be a $K$-vector space with ordered bases $\mathcal{B}$ and $\mathcal{C}$.
    \begin{enumerate}
        \item \textbf{Coordinate conversion}:
              $\CoordVec[\mathcal{C}]{v}
              = \CoordMatrix[\mathcal{B}][\mathcal{C}]{\Identity{V}}
                \cdot \CoordVec[\mathcal{B}]{v}$
              for all $v \in V$.
        \item \textbf{Invertibility}:
              $\CoordMatrix[\mathcal{B}][\mathcal{C}]{\Identity{V}}
               \cdot \CoordMatrix[\mathcal{C}][\mathcal{B}]{\Identity{V}} = I$,
               i.e.\
              $\bigl(\CoordMatrix[\mathcal{B}][\mathcal{C}]{\Identity{V}}\bigr)^{-1}
              = \CoordMatrix[\mathcal{C}][\mathcal{B}]{\Identity{V}}$.
        \item The change of basis formula (Theorem~\ref{thm:change-of-bases}) reads:
              \[
                  \CoordMatrix[\mathcal{B}_U][\mathcal{B}_V]{T}
                  = \CoordMatrix[\mathcal{C}_V][\mathcal{B}_V]{\Identity{V}}
                    \cdot \CoordMatrix[\mathcal{C}_U][\mathcal{C}_V]{T}
                    \cdot \CoordMatrix[\mathcal{B}_U][\mathcal{C}_U]{\Identity{U}}
              \]
    \end{enumerate}
\end{proposition}

\begin{proof}
    (1) is the coordinate formula (Corollary~\ref{prop:coord-formula}) applied to
    $\Identity{V}$.
    (2) follows from Corollary~\ref{thm:comp-mult} applied to
    $\Composition{\Identity{V}}{\Identity{V}} = \Identity{V}$:
    \[
        \CoordMatrix[\mathcal{B}][\mathcal{C}]{\Identity{V}}
        \cdot \CoordMatrix[\mathcal{C}][\mathcal{B}]{\Identity{V}}
        = \CoordMatrix[\mathcal{B}][\mathcal{B}]{\Identity{V}}
        = I \qedhere
    \]
\end{proof}

\subsection{Change of Basis as an Automorphism of $K^n$}

\begin{remark}[Automorphism Interpretation]
\label{rem:cob-automorphism}
    When $\Dim[K]{V} = n$, each ordered basis $\mathcal{B}$ identifies $V \cong K^n$ via
    $\phi_{\mathcal{B}}$. A second basis $\mathcal{C}$ gives a second identification
    $V \cong K^n$ via $\phi_{\mathcal{C}}$. Both land in the \emph{same} $K^n$, encoding each
    $v \in V$ as two different coordinate tuples:
    $\CoordVec[\mathcal{B}]{v} \in K^n$ and $\CoordVec[\mathcal{C}]{v} \in K^n$.

    The change of basis matrix
    $\CoordMatrix[\mathcal{B}][\mathcal{C}]{\Identity{V}} \in \CFMat{n}{n}[K]$,
    viewed as a linear map $K^n \to K^n$ (Section~\ref{sec:mat-as-lin-maps}), is an
    \textbf{automorphism of $K^n$}: it is an invertible endomorphism that converts
    $\mathcal{B}$-coordinates into $\mathcal{C}$-coordinates for every vector simultaneously.
    This is the correct setting in which to speak of ``change of basis'': it is not a change
    of the space $K^n$, but a change of how vectors in the same space are coordinatised.
\end{remark}

\subsection{Examples}

\begin{example}[Computing a Change of Basis Matrix]
\label{ex:cob-R2}
    Let $V = \bb{R}^2$, $\mathcal{B} = (e_0, e_1)$ (standard),
    $\mathcal{C} = (e_0 + e_1,\, e_1)$.

    Express each $b_j \in \mathcal{B}$ in $\mathcal{C}$:
    $e_0 = 1 \cdot (e_0+e_1) + (-1) \cdot e_1$, so
    $\CoordVec[\mathcal{C}]{e_0} = \begin{psmallmatrix}1\\-1\end{psmallmatrix}$; and
    $e_1 = 0 \cdot (e_0+e_1) + 1 \cdot e_1$, so
    $\CoordVec[\mathcal{C}]{e_1} = \begin{psmallmatrix}0\\1\end{psmallmatrix}$.
    Hence:
    \[
        \CoordMatrix[\mathcal{B}][\mathcal{C}]{\Identity{}}
        = \begin{pmatrix} 1 & 0 \\ -1 & 1 \end{pmatrix}
    \]
    Verify coordinate conversion: for $v = (3,5)$,
    $\CoordVec[\mathcal{B}]{v} = \begin{psmallmatrix}3\\5\end{psmallmatrix}$ and
    \[
        \CoordMatrix[\mathcal{B}][\mathcal{C}]{\Identity{}}
        \cdot \CoordVec[\mathcal{B}]{v}
        = \begin{pmatrix}1&0\\-1&1\end{pmatrix}\begin{pmatrix}3\\5\end{pmatrix}
        = \begin{pmatrix}3\\2\end{pmatrix}
        = \CoordVec[\mathcal{C}]{v}
    \]
    Direct check: $v = (3,5) = 3(e_0+e_1) + 2 e_1$ $\checkmark$.

    The inverse:
    $\CoordMatrix[\mathcal{C}][\mathcal{B}]{\Identity{}}
    = \begin{pmatrix}1&0\\1&1\end{pmatrix}$ (columns are $\mathcal{C}$-vectors in
    $\mathcal{B}$-coords), and indeed
    $\begin{pmatrix}1&0\\-1&1\end{pmatrix}
     \begin{pmatrix}1&0\\1&1\end{pmatrix}
     = I$ $\checkmark$.
\end{example}

\begin{example}[Change of Basis for an Endomorphism]
    With $T(x,y) = (2x+y,\, x)$ and bases $\mathcal{B}$, $\mathcal{C}$ from
    Example~\ref{ex:cob-R2}, the change of bases formula gives (cf.\
    Example~\ref{ex:cob-endomorphism}):
    \[
        \CoordMatrix[\mathcal{C}][\mathcal{C}]{T}
        = \CoordMatrix[\mathcal{B}][\mathcal{C}]{\Identity{}}
          \cdot \CoordMatrix[\mathcal{B}][\mathcal{B}]{T}
          \cdot \CoordMatrix[\mathcal{C}][\mathcal{B}]{\Identity{}}
        = \begin{pmatrix}1&0\\-1&1\end{pmatrix}
          \begin{pmatrix}2&1\\1&0\end{pmatrix}
          \begin{pmatrix}1&0\\1&1\end{pmatrix}
        = \begin{pmatrix} 3 & 1 \\ -2 & -1 \end{pmatrix}
    \]
\end{example}

\subsection{Injectivity, Surjectivity, Bijectivity via Matrices}

\begin{remark}[Substitution Principle]
\label{rem:substitution-principle}
    Every characterization of $T$ from the theorems of
    Section~\ref{sec:inj-surj-bij} yields an equivalent characterization with $T$ replaced
    by any matrix $\CoordMatrix[\mathcal{B}_U][\mathcal{B}_V]{T}$ representing it.
    Precisely: $T$ satisfies property $P$ if and only if
    $\CoordMatrix[\mathcal{B}_U][\mathcal{B}_V]{T}$ satisfies property $P$.

    \textit{Reason}: $\CoordMatrix[\mathcal{B}_U][\mathcal{B}_V]{T}
    = \Composition{\phi_{\mathcal{B}_V}}{\Composition{T}{\phi_{\mathcal{B}_U}^{-1}}}$
    is a conjugate of $T$ by isomorphisms. Isomorphisms preserve and reflect every structural
    property: injectivity, surjectivity, kernel dimension, image dimension, etc.
\end{remark}

\begin{theorem}[Matrix Characterizations of Injectivity]
\label{thm:mat-inj}
    Let $T: U \to V$ be $K$-linear with $\kappa_1 = \Dim[K]{U} \leq \Dim[K]{V} = \kappa_2$.
    Each of the following is equivalent to $T$ being injective:
    \begin{enumerate}
        \item For every ordered basis $\mathcal{B}_U$ of $U$, there exists an ordered basis
              $\mathcal{B}_V$ of $V$ such that:
              \[
                  \CoordMatrix[\mathcal{B}_U][\mathcal{B}_V]{T}
                  = \begin{pmatrix} I_{\kappa_1} \\ 0 \end{pmatrix}
              \]
        \item There exist ordered bases $\mathcal{B}_U$ of $U$ and $\mathcal{B}_V$ of $V$
              such that:
              \[
                  \CoordMatrix[\mathcal{B}_U][\mathcal{B}_V]{T}
                  = \begin{pmatrix} I_{\kappa_1} \\ 0 \end{pmatrix}
              \]
        \item For every $w \in V$, the system
              $\CoordMatrix[\mathcal{B}_U][\mathcal{B}_V]{T}\, x = \CoordVec[\mathcal{B}_V]{w}$
              has at most one solution $x \in K^{(\kappa_1)}$ (for any ordered bases).
        \item The columns of $\CoordMatrix[\mathcal{B}_U][\mathcal{B}_V]{T}$ are linearly
              independent (for any ordered bases $\mathcal{B}_U$, $\mathcal{B}_V$).
    \end{enumerate}
\end{theorem}

\begin{proof}
    \leavevmode
    \begin{enumerate}
        \item Reformulation of: ``$T$ sends every basis to a linearly independent set.''
              Given $\mathcal{B}_U = (b_1, \ldots, b_{\kappa_1})$, by injectivity
              $\{T(b_1), \ldots, T(b_{\kappa_1})\}$ is linearly independent in $V$. Extend
              it to an ordered basis
              $\mathcal{B}_V = (T(b_1), \ldots, T(b_{\kappa_1}),
              c_{\kappa_1+1}, \ldots, c_{\kappa_2})$ of $V$.
              The $j$-th column of $\CoordMatrix[\mathcal{B}_U][\mathcal{B}_V]{T}$ is
              $\CoordVec[\mathcal{B}_V]{T(b_j)} = e_j \in K^{(\kappa_2)}$, giving exactly
              $\begin{pmatrix} I_{\kappa_1} \\ 0 \end{pmatrix}$. Conversely, if such a
              $\mathcal{B}_V$ exists then the columns are the first $\kappa_1$ standard
              basis vectors, hence linearly independent, so $T$ is injective
              (Remark~\ref{rem:substitution-principle}).

        \item Reformulation of: ``$T$ sends some basis to a linearly independent set.''
              The same argument applied to one chosen basis.

        \item Reformulation of set-theoretic injectivity: the system
              $\CoordMatrix[\mathcal{B}_U][\mathcal{B}_V]{T}\, x
              = \CoordVec[\mathcal{B}_V]{w}$
              is the coordinate form of $T(v) = w$ (coordinate formula). At most one
              coordinate solution $x$ is equivalent to at most one preimage
              $v = \phi_{\mathcal{B}_U}^{-1}(x)$.

        \item Reformulation of: ``$T$ sends every basis to a linearly independent set.''
              The $j$-th column of $\CoordMatrix[\mathcal{B}_U][\mathcal{B}_V]{T}$ is
              $\phi_{\mathcal{B}_V}(T(b_j))$.
              Since $\phi_{\mathcal{B}_V}$ is an isomorphism, the columns are linearly
              independent iff $T(b_1), \ldots, T(b_{\kappa_1})$ are linearly independent
              iff $T$ sends $\mathcal{B}_U$ to a linearly independent set.
              This holds for every $\mathcal{B}_U$ iff $T$ sends every basis to a linearly
              independent set, which is characterization~(7) of
              Theorem~\ref{thm:injectivity}. \qedhere
    \end{enumerate}
\end{proof}

\begin{theorem}[Matrix Characterizations of Surjectivity]
\label{thm:mat-surj}
    Let $T: U \to V$ be $K$-linear with $\kappa_2 = \Dim[K]{V} \leq \Dim[K]{U} = \kappa_1$.
    Each of the following is equivalent to $T$ being surjective:
    \begin{enumerate}
        \item For every ordered basis $\mathcal{B}_V$ of $V$, there exists an ordered basis
              $\mathcal{B}_U$ of $U$ such that:
              \[
                  \CoordMatrix[\mathcal{B}_U][\mathcal{B}_V]{T}
                  = \begin{pmatrix} I_{\kappa_2} & 0 \end{pmatrix}
              \]
        \item There exist ordered bases $\mathcal{B}_U$ of $U$ and $\mathcal{B}_V$ of $V$
              such that:
              \[
                  \CoordMatrix[\mathcal{B}_U][\mathcal{B}_V]{T}
                  = \begin{pmatrix} I_{\kappa_2} & 0 \end{pmatrix}
              \]
        \item For every $w \in V$, the system
              $\CoordMatrix[\mathcal{B}_U][\mathcal{B}_V]{T}\, x = \CoordVec[\mathcal{B}_V]{w}$
              has at least one solution $x \in K^{(\kappa_1)}$ (for any ordered bases).
        \item The columns of $\CoordMatrix[\mathcal{B}_U][\mathcal{B}_V]{T}$ generate
              $K^{(\kappa_2)}$ (for any ordered bases $\mathcal{B}_U$, $\mathcal{B}_V$).
    \end{enumerate}
\end{theorem}

\begin{proof}
    \leavevmode
    \begin{enumerate}
        \item Reformulation of: ``$T$ sends every basis to a generating set.'' Given
              $\mathcal{B}_V = (c_1, \ldots, c_{\kappa_2})$, since $T$ is surjective pick
              $u_i \in U$ with $T(u_i) = c_i$ for each $i$. The set
              $\{u_1, \ldots, u_{\kappa_2}\}$ is linearly independent (as $T$ is injective
              on $\Span[K]{\{u_1,\ldots,u_{\kappa_2}\}}$ via the coordinate argument).
              Extend by choosing $u_{\kappa_2+1}, \ldots, u_{\kappa_1} \in \Ker{T}$ forming
              a basis of $\Ker{T}$ (which has dimension $\kappa_1 - \kappa_2$ by
              rank-nullity), giving basis $\mathcal{B}_U = (u_1, \ldots, u_{\kappa_1})$.
              The first $\kappa_2$ columns of $\CoordMatrix[\mathcal{B}_U][\mathcal{B}_V]{T}$
              are $e_1, \ldots, e_{\kappa_2}$ (since $T(u_i) = c_i$) and the last
              $\kappa_1 - \kappa_2$ columns are zero (since $T(u_j) = 0$ for
              $j > \kappa_2$), giving $\begin{pmatrix} I_{\kappa_2} & 0 \end{pmatrix}$.

        \item Reformulation of: ``$T$ sends some basis to a generating set.'' Same argument
              for one chosen $\mathcal{B}_V$.

        \item Reformulation of surjectivity: every $w \in V$ has at least one preimage, i.e.\
              the system $\CoordMatrix[\mathcal{B}_U][\mathcal{B}_V]{T}\, x
              = \CoordVec[\mathcal{B}_V]{w}$ always has a solution.

        \item Reformulation of: ``$T$ sends every basis to a generating set.''
              The $j$-th column of $\CoordMatrix[\mathcal{B}_U][\mathcal{B}_V]{T}$
              is $\phi_{\mathcal{B}_V}(T(b_j))$.
              Since $\phi_{\mathcal{B}_V}$ is an isomorphism, the columns generate
              $K^{(\kappa_2)}$ iff $T(b_1), \ldots, T(b_{\kappa_1})$ generate $V$ iff $T$
              sends $\mathcal{B}_U$ to a generating set. This holds for every $\mathcal{B}_U$
              iff $T$ sends every basis to a generating set, which is
              characterization~(7) of Theorem~\ref{thm:surjectivity}. \qedhere
    \end{enumerate}
\end{proof}

\begin{theorem}[Matrix Characterizations of Bijectivity]
\label{thm:mat-bij}
    Let $T: U \to V$ be $K$-linear with $\kappa_1 = \Dim[K]{U} = \Dim[K]{V} = \kappa_2$.
    Each of the following is equivalent to $T$ being bijective:
    \begin{enumerate}
        \item For every ordered basis $\mathcal{B}_U$ of $U$, there exists an ordered basis
              $\mathcal{B}_V$ of $V$ such that $\CoordMatrix[\mathcal{B}_U][\mathcal{B}_V]{T} = I_{\kappa_1}$.
        \item For every ordered basis $\mathcal{B}_V$ of $V$, there exists an ordered basis
              $\mathcal{B}_U$ of $U$ such that $\CoordMatrix[\mathcal{B}_U][\mathcal{B}_V]{T} = I_{\kappa_1}$.
        \item There exist ordered bases $\mathcal{B}_U$ of $U$ and $\mathcal{B}_V$ of $V$
              such that $\CoordMatrix[\mathcal{B}_U][\mathcal{B}_V]{T} = I_{\kappa_1}$.
        \item For every $w \in V$, the system
              $\CoordMatrix[\mathcal{B}_U][\mathcal{B}_V]{T}\, x = \CoordVec[\mathcal{B}_V]{w}$
              has exactly one solution $x \in K^{(\kappa_1)}$ (for any ordered bases).
        \item The columns of $\CoordMatrix[\mathcal{B}_U][\mathcal{B}_V]{T}$ form a basis of
              $K^{(\kappa_1)}$ (for any ordered bases $\mathcal{B}_U$, $\mathcal{B}_V$).
    \end{enumerate}
\end{theorem}

\begin{proof}
    \leavevmode
    \begin{enumerate}
        \item Reformulation of: ``$T$ sends every basis to a basis''
              (Theorem~\ref{thm:bijectivity}). Given
              $\mathcal{B}_U = (b_1, \ldots, b_{\kappa_1})$, by bijectivity
              $T(\mathcal{B}_U)$ is a basis of $V$. Set
              $\mathcal{B}_V = (T(b_1), \ldots, T(b_{\kappa_1}))$; then the $j$-th column
              is $\CoordVec[\mathcal{B}_V]{T(b_j)} = e_j$, so
              $\CoordMatrix[\mathcal{B}_U][\mathcal{B}_V]{T} = I_{\kappa_1}$.

        \item Given $\mathcal{B}_V = (c_1, \ldots, c_{\kappa_1})$, since $T$ is bijective,
              the preimages $b_j = T^{-1}(c_j)$ form a basis $\mathcal{B}_U$ of $U$. The
              $j$-th column is $\CoordVec[\mathcal{B}_V]{T(b_j)} = \CoordVec[\mathcal{B}_V]{c_j}
              = e_j$, so $\CoordMatrix[\mathcal{B}_U][\mathcal{B}_V]{T} = I_{\kappa_1}$.

        \item Follows from (1) or (2).

        \item Combine Theorem~\ref{thm:mat-inj}(3) (at most one solution) and
              Theorem~\ref{thm:mat-surj}(3) (at least one solution).

        \item Reformulation of: ``$T$ sends every basis to a basis.''
              The columns are $\phi_{\mathcal{B}_V}(T(b_j))$; since $\phi_{\mathcal{B}_V}$
              is an isomorphism, they form a basis of $K^{(\kappa_1)}$ iff $T(\mathcal{B}_U)$
              is a basis of $V$, which holds for every $\mathcal{B}_U$ iff $T$ sends every
              basis to a basis --- characterization~(5) of
              Theorem~\ref{thm:bijectivity}. \qedhere
    \end{enumerate}
\end{proof}

\begin{corollary}[Automorphisms Are Exactly the Change of Basis Matrices]
\label{cor:aut-cob}
    Let $A \in \CFMat{\kappa}{\kappa}[K]$, viewed as an endomorphism of $K^{(\kappa)}$.
    The following are equivalent:
    \begin{enumerate}
        \item $A$ is an automorphism of $K^{(\kappa)}$.
        \item For every ordered basis $\mathcal{B}$ of $K^{(\kappa)}$, there exists an ordered
              basis $\mathcal{C}$ of $K^{(\kappa)}$ such that
              $A = \CoordMatrix[\mathcal{B}][\mathcal{C}]{\Identity{K^{(\kappa)}}}$.
        \item For every ordered basis $\mathcal{C}$ of $K^{(\kappa)}$, there exists an ordered
              basis $\mathcal{B}$ of $K^{(\kappa)}$ such that
              $A = \CoordMatrix[\mathcal{B}][\mathcal{C}]{\Identity{K^{(\kappa)}}}$.
        \item There exist ordered bases $\mathcal{B}$, $\mathcal{C}$ of $K^{(\kappa)}$ such
              that $A = \CoordMatrix[\mathcal{B}][\mathcal{C}]{\Identity{K^{(\kappa)}}}$.
    \end{enumerate}
\end{corollary}

\begin{proof}
    Let $B = [b_1 \mid \cdots \mid b_\kappa]$ and $C = [c_1 \mid \cdots \mid c_\kappa]$
    denote the matrices whose columns are the basis vectors of $\mathcal{B}$ and
    $\mathcal{C}$. The $j$-th column of
    $\CoordMatrix[\mathcal{B}][\mathcal{C}]{\Identity{}}$ is $\CoordVec[\mathcal{C}]{b_j}$,
    i.e.\ the unique $x$ with $Cx = b_j$, so
    $\CoordMatrix[\mathcal{B}][\mathcal{C}]{\Identity{}} = C^{-1}B$.
    Thus $A = C^{-1}B$ iff $B = CA$.

    \begin{itemize}
        \item (1)$\Rightarrow$(2): Given $\mathcal{B}$ (with matrix $B$), set $C = BA^{-1}$.
              Since $B$ and $A^{-1}$ are invertible, $C$ is invertible, so $\mathcal{C}$
              (columns of $C$) is a basis. Then
              $\CoordMatrix[\mathcal{B}][\mathcal{C}]{\Identity{}} = C^{-1}B
              = (BA^{-1})^{-1}B = AB^{-1}B = A$.

        \item (1)$\Rightarrow$(3): Given $\mathcal{C}$ (with matrix $C$), set $B = CA$.
              Since $C$ and $A$ are invertible, $B$ is invertible, so $\mathcal{B}$
              (columns of $B$) is a basis. Then
              $\CoordMatrix[\mathcal{B}][\mathcal{C}]{\Identity{}} = C^{-1}B = C^{-1}CA = A$.

        \item (2)$\Rightarrow$(4) and (3)$\Rightarrow$(4): Immediate.

        \item (4)$\Rightarrow$(1): Every change of basis matrix is invertible
              (Proposition~\ref{prop:cob-properties}(2)), hence an automorphism. \qedhere
    \end{itemize}
\end{proof}

\begin{example}[Normal Form for Injection]
    $T: \bb{R}^2 \to \bb{R}^3$, $T(x,y) = (x, y, 0)$.
    With $\mathcal{B}_U = (e_0, e_1)$ and $\mathcal{B}_V = (e_0, e_1, e_2)$ (standard), already
    $\CoordMatrix[\mathcal{B}_U][\mathcal{B}_V]{T}
    = \begin{psmallmatrix}1&0\\0&1\\0&0\end{psmallmatrix}
    = \begin{psmallmatrix}I_2\\0\end{psmallmatrix}$.
\end{example}

\begin{example}[Normal Form for Surjection]
    $T: \bb{R}^3 \to \bb{R}^2$, $T(x,y,z) = (x,y)$.
    Take $\mathcal{B}_U = (e_0, e_1, e_2)$, $\mathcal{B}_V = (e_0, e_1)$:
    $\CoordMatrix[\mathcal{B}_U][\mathcal{B}_V]{T}
    = \begin{pmatrix}1&0&0\\0&1&0\end{pmatrix}
    = \begin{pmatrix}I_2 & 0\end{pmatrix}$.
\end{example}

\begin{example}[No Normal Form]
    $T: \bb{R}^2 \to \bb{R}^2$, $T(x,y) = (x+y, x+y)$. Since
    $\Rank{T} = 1 < 2$, no choice of bases gives $I_2$; $T$ is neither injective
    nor surjective.
\end{example}

\subsection{Invertibility}

\begin{theorem}[Inverse via Matrix Inverse]
\label{thm:mat-inverse}
    Let $T: U \xrightarrow{\sim} V$ be a $K$-linear isomorphism with ordered bases
    $\mathcal{B}_U$, $\mathcal{B}_V$. Then:
    \[
        \bigl(\CoordMatrix[\mathcal{B}_U][\mathcal{B}_V]{T}\bigr)^{-1}
        = \CoordMatrix[\mathcal{B}_V][\mathcal{B}_U]{T^{-1}}
    \]
\end{theorem}

\begin{proof}
    Apply Corollary~\ref{thm:comp-mult} to
    $\Composition{T^{-1}}{T} = \Identity{U}$ and
    $\Composition{T}{T^{-1}} = \Identity{V}$:
    \begin{align*}
        \CoordMatrix[\mathcal{B}_V][\mathcal{B}_U]{T^{-1}}
        \cdot \CoordMatrix[\mathcal{B}_U][\mathcal{B}_V]{T}
        &= \CoordMatrix[\mathcal{B}_U][\mathcal{B}_U]{\Composition{T^{-1}}{T}}
         = \CoordMatrix[\mathcal{B}_U][\mathcal{B}_U]{\Identity{U}}
         = I_n \\
        \CoordMatrix[\mathcal{B}_U][\mathcal{B}_V]{T}
        \cdot \CoordMatrix[\mathcal{B}_V][\mathcal{B}_U]{T^{-1}}
        &= \CoordMatrix[\mathcal{B}_V][\mathcal{B}_V]{\Composition{T}{T^{-1}}}
         = \CoordMatrix[\mathcal{B}_V][\mathcal{B}_V]{\Identity{V}}
         = I_m
    \end{align*}
    Hence the two matrices are mutual inverses. \qedhere
\end{proof}

\begin{corollary}[$T$ Invertible $\Leftrightarrow$ Matrix Invertible]
\label{cor:invertible-iff-mat-invertible}
    $T: U \to V$ is an isomorphism if and only if
    $\CoordMatrix[\mathcal{B}_U][\mathcal{B}_V]{T}$ is an invertible matrix (for any,
    equivalently some, ordered bases $\mathcal{B}_U$, $\mathcal{B}_V$).
\end{corollary}

\begin{proof}
    ($\Rightarrow$) Theorem~\ref{thm:mat-inverse} gives an explicit matrix inverse.

    ($\Leftarrow$) Let $A = \CoordMatrix[\mathcal{B}_U][\mathcal{B}_V]{T}$ be invertible.
    Define $S \defeq \Composition{\phi_{\mathcal{B}_U}^{-1}}
    {\Composition{A^{-1}}{\phi_{\mathcal{B}_V}}}: V \to U$.
    Then $S$ is $K$-linear and:
    \[
        \CoordMatrix[\mathcal{B}_V][\mathcal{B}_U]{S}
        \cdot \CoordMatrix[\mathcal{B}_U][\mathcal{B}_V]{T}
        = A^{-1} \cdot A = I_n
        = \CoordMatrix[\mathcal{B}_U][\mathcal{B}_U]{\Identity{U}}
    \]
    so $\Composition{S}{T} = \Identity{U}$; similarly $\Composition{T}{S} = \Identity{V}$.
    Hence $T$ is an isomorphism.

    \textit{Independence of bases}: If $A = \CoordMatrix[\mathcal{B}_U][\mathcal{B}_V]{T}$
    is invertible, then for any other bases $\mathcal{C}_U$, $\mathcal{C}_V$:
    \[
        \CoordMatrix[\mathcal{C}_U][\mathcal{C}_V]{T}
        = \CoordMatrix[\mathcal{B}_V][\mathcal{C}_V]{\Identity{V}}
          \cdot A
          \cdot \CoordMatrix[\mathcal{C}_U][\mathcal{B}_U]{\Identity{U}}
    \]
    is a product of invertible matrices, hence invertible. \qedhere
\end{proof}

\subsection{Examples}

\begin{example}[Inverse via Matrix Inverse]
\label{ex:mat-inverse}
    Let $T: \bb{R}^2 \to \bb{R}^2$, $T(x,y) = (2x+y,\, x+y)$, with standard bases.
    \[
        A = \CoordMatrix[\mathcal{E}][\mathcal{E}]{T}
        = \begin{pmatrix} 2 & 1 \\ 1 & 1 \end{pmatrix},
        \qquad
        A^{-1} = \begin{pmatrix} 1 & -1 \\ -1 & 2 \end{pmatrix}
    \]
    By Theorem~\ref{thm:mat-inverse},
    $\CoordMatrix[\mathcal{E}][\mathcal{E}]{T^{-1}} = A^{-1}$, so
    $T^{-1}(x,y) = (x-y,\, -x+2y)$.
    Verify: $T(T^{-1}(x,y)) = T(x-y,\, -x+2y) = (2(x-y)+(-x+2y),\, (x-y)+(-x+2y))
    = (x, y)$ $\checkmark$.
\end{example}

\begin{example}[Not Invertible]
    $T: \bb{R}^2 \to \bb{R}^2$, $T(x,y) = (x+y,\, x+y)$.
    $\CoordMatrix[\mathcal{E}][\mathcal{E}]{T}
    = \begin{pmatrix}1&1\\1&1\end{pmatrix}$: the two columns are equal, hence linearly
    dependent, so $T$ is not injective and not an isomorphism
    (Theorem~\ref{thm:mat-inj}(4)).
\end{example}

\begin{example}[Basis-Independence of Invertibility]
    For $T$ from Example~\ref{ex:mat-inverse} with $\mathcal{C} = (e_0+e_1,\, e_1)$:
    \[
        \CoordMatrix[\mathcal{C}][\mathcal{C}]{T}
        = \begin{pmatrix}3&1\\-2&-1\end{pmatrix}
    \]
    The columns $\begin{psmallmatrix}3\\-2\end{psmallmatrix}$ and
    $\begin{psmallmatrix}1\\-1\end{psmallmatrix}$ are linearly independent (neither is a
    scalar multiple of the other), so $T$ is injective, and since $n = m = 2$ it is an
    isomorphism --- as expected from Corollary~\ref{cor:invertible-iff-mat-invertible}.
\end{example}
